\section{Introdução}

\hspace{0.4cm}Apresentando-se como um método construtivo que possibilita a elaboração de grandes vãos, Veríssimo e Junior (\citeyear{VER&JUN}) entendem que a consagração das estruturas de concreto protendido (CP) se dá através do grande número de obras diversificadas realizadas. Contudo como explica Bastos (\citeyear{BAS}), o CP inicialmente não teve sucesso devido à baixa resistência dos aços de protensão e às perdas significativas de força ao longo do tempo. Foi somente durante a Segunda Guerra Mundial, com o desenvolvimento de aços de alta resistência, que a técnica se tornou viável, impulsionada pelo engenheiro francês Eugène Freyssinet. Após a guerra, devido à necessidade de reconstrução de infraestruturas na Europa, a protensão ganhou ainda mais notoriedade. Nos Estados Unidos surgiram as primeiras patentes de protensão, enquanto se desenvolveu o conceito de protensão parcial na Inglaterra.

Entende-se que o método se expandiu globalmente, com contribuições significativas de engenheiros em vários países. Ele é descrito por Pfeil (\citeyear{PFEA}) como uma técnica que envolve adicionar tensões previamente em uma estrutura para melhorar sua resistência e desempenho sob diferentes condições de carga. Carvalho (\citeyear{CAR}) analisa que no Brasil, as estruturas de CP e de concreto armado (CA) são encaradas como do mesmo tipo, visto que são regidas pelo mesmo documento no país, e dado que suas principais divergências estão no tipo de aço empregado e processo construtivo, utiliza de diferentes especificações quando é o caso. O documento descrito trata-se da Norma Brasileira (NBR) 6118, que sofre atualizações esporádicas, adquirindo em sua nomenclatura o ano delas.

Para o desenvolvimento de projetos estruturais de CP, necessita-se da aplicação de uma série de técnicas e cálculos, que particularmente se distinguem dos utilizados no dimensionamento de CA, com especial atenção ao controle da força de protensão, que garante o desempenho adequado da estrutura. Dentre eles destaca-se o abordado no presente estudo, que é o calculo das perdas de protensão. Cholfe e Bonilha (\citeyear{CB}) entendem que em um projeto de CP, o elemento fundamental é a força de protensão, que depende dos componentes físicos das ferramentas de protensão e do comportamento intrínseco dos materiais utilizados, e ela relacionada ao aparelho tensor sofre de perdas que devem ser previstas. Estas ocorrem antes da transferência da protensão ao concreto como perdas iniciais, durante a transferência como perdas imediatas, e também ao longo da vida útil da estrutura como perdas progressivas.

Contudo, muitas vezes essa etapa do dimensionamento de estruturas de CP é negligenciada pelos projetistas, ao passo que autores como Pfeil (\citeyear{PFEA}) e Boukendakdji \textit{et al.} (\citeyear{BEA}) declaram que é comum o uso de constantes para as perdas de protensão de forma que os projetos acabam superdimensionados, elevando de forma desnecessária o custo da estrutura. Além disso, existem ferramentas computacionais pagas que se propõem a fazer o cálculo estrutural completo de estruturas de CP. 

Porém dado a automatização e complexidade envolvida nos cálculos realizados por esse tipo de programa, etapas como a da perda de protensão acabam por ficar em segundo plano. Volta-se então ao prática já citada na definição das perdas de protensão por meio de constantes e valores definidos pelo usuário, a exemplo do TQS com as perdas progressivas de protensão, em que seu valor deve ser fornecido em porcentagem (\citeyear{TQA}). Revelam-se assim, sistemas não transparente em meio aos cálculos estruturais das ferramentas, dado sua baixa explicabilidade, tornando assim essencial o uso de alternativas auxiliares para validação e aumento da clareza do dimensionamento

Ainda sim, existe a possibilidade da realização dos cálculos de forma manual. Entretanto, Mascarenhas \textit{et al.} (\citeyear{MEA}) julgam que o grande número de equações, processos e complexidade dos cálculos podem exigir uma grande parcela tempo que irá compor o desenvolvimento do projeto. Seguindo essa ideia, torna-se uma inviável o cálculo manual, dado que o prazo de entrega pode acabar não sendo competitivo em comparação a outros métodos construtivos. Labadan (\citeyear{LAB}) complementa esse argumento, ao colocar que a natureza iterativa dos cálculos de estruturas de CP para que se chegue ao Estado Limite de Serviço (ELS) e ao Estado Limite Último (ELU), podem levar projetistas que realizem o dimensionamento de forma manual dado a impaciência, assumindo valores de forma precipitada e chegando assim a proporções não econômicas.

\newpage

Assim sendo, com o objetivo de elevar a clareza nos cálculos que de perda de protensão, de maneira que aumente a capacidade de modificação e explicação deles, fomentando o dimensionamento de estruturas mais bem otimizadas, e por consequência diminuindo seu custo, o presente estudo se propõe a desenvolver uma ferramenta computacional acessível, gratuita e de código aberto, para a realização dos cálculos de perda de protensão. Ela será desenvolvida através da linguagem de programação \textit{Python}, com um instalador e interface gráfica interativa e intuitiva para os projetistas.